% Generated draft: requires mathematical and source review.
\documentclass[11pt]{article}
\usepackage[margin=25mm]{geometry}
\usepackage{fontspec}
\setmainfont{Latin Modern Roman}
\usepackage{amsmath,amssymb,mathtools,unicode-math}
\setmathfont{Latin Modern Math}
\usepackage{enumitem,needspace,xurl,hyperref}
\setlength{\parindent}{0pt}
\setlength{\parskip}{5pt}
\setlength{\emergencystretch}{3em}
\newcommand{\N}{\mathbb{N}}
\newcommand{\Z}{\mathbb{Z}}
\newcommand{\Q}{\mathbb{Q}}
\newcommand{\R}{\mathbb{R}}
\newcommand{\C}{\mathbb{C}}
\newcommand{\F}{\mathbb{F}}
\newcommand{\A}{\mathbb{A}}
\newcommand{\PP}{\mathbb{P}}
\newcommand{\E}{\mathbb{E}}
\newcommand{\eps}{\varepsilon}
\newcommand{\dd}{\,\mathrm{d}}
\begin{document}
\section*{UnsolvedMath ID: 20003157}
\textbf{Source problem code:} catalogue entry
\subsection{Definitions and mathematical statement}
The source asks: if a group acts isometrically on a manifold with \(\sec\geq1\), and its entire fixed-point set is a circle, must the induced length be at most \(2\pi\)? The extracted trailing heading ``4. Manifolds of Cohomogeneity-one and Polar Actions'' belongs to the next section, not the mathematical hypothesis. Completeness, dimension, compactness of the acting group, and whether one means the whole fixed set or a component are not otherwise specified.

The aid's literal counterexample uses \(R>1\), \(0<a<\pi/2\), and the boundaryless open cylinder
\[
 M_{R,a}=(-a,a)\times(\R/2\pi\Z),\qquad
 g=\dd r^2+R^2\cos^2r\,\dd\theta^2.
\]
Its Gaussian curvature is identically \(1\). The effective \(\Z/2\Z\)-action generated by
\[
 \iota(r,\theta)=(-r,\theta)
\]
is isometric, and
\[
 \operatorname{Fix}(\iota)=\{0\}\times(\R/2\pi\Z),\qquad
 \operatorname{length}\operatorname{Fix}(\iota)=2\pi R>2\pi.
\]
The missing ends are at finite radial distance, so the metric is incomplete. Varying \(R\) makes the fixed-circle length arbitrarily large while keeping curvature one. This disproves the literal assertion without completeness and does not disprove its complete version.

For the likely intended complete version in dimensions at least two, the aid separately reports the sharp bound in dimension two. It also reports the bound in higher dimensions under a slice-representation condition: if \(C=\operatorname{Fix}(G)\) is the circle and, at a point \(p\in C\), some subgroup of the isotropy action fixes exactly one normal direction, meaning
\[
 \dim(\nu_pC)^H=1\quad\text{for a subgroup }H\leq G,
\]
then \(\operatorname{length}(C)\leq2\pi\), under the complete geometric hypotheses used in that argument. The detailed acting-group and fixed-component assumptions of this partial bound are not supplied. The standard unit sphere with equatorial reflection exhibits equality in dimension two.

The further Morse-theoretic observation reported by the aid is that for any closed geodesic \(\gamma\) of length \(L>2\pi\) in an \(n\)-dimensional manifold with \(\sec\geq1\),
\[
 \operatorname{index}(\gamma)\geq2(n-1),
\]
where index counts negative directions of the closed-geodesic second variation. This index statement does not itself exclude a fixed geodesic in every group action. The source supplies no general resolution of the complete higher-dimensional fixed-circle question outside its stated sufficient slice condition. The registry's solved label is therefore retained only as the aid's assessment of the incomplete literal formulation.
\subsection{Short English statement}
The supplied incomplete constant-curvature cylinder has a reflected fixed circle of arbitrary length, so completeness matters. Separate partial bounds concern complete surfaces and a stated normal-slice condition; they are not a solution of every complete higher-dimensional case.
\subsection{Sources}
\begin{itemize}
\item Local UnsolvedMath registry, permanent ID \(20003157\), exact source code \texttt{catalogue entry}; canonical \texttt{aim-topology-notes.json}, \texttt{notes[244]}, source tag \texttt{problem}.
\item Problem 32 in M. Kerin's compilation from the 2007 AIM workshop ``Manifolds with Non-negative Sectional Curvature'': \url{https://aimath.org/WWN/nnsectcurvature/nnsectcurvature.pdf}. The supplied verification calls the canonical wording exact. The next section heading is extraction spillover. No PDF was fetched.
\item The original literature field is empty. No complete authoritative citation for the complete-case partial bounds or the Morse-index statement is supplied; the explicit counterexample and partial results are reported by the research aid.
\item Research attempt 1: result \texttt{counterexample}, problem status at run \texttt{unknown}, rigor label \texttt{proved}, novelty confidence \texttt{low}; classification \texttt{SOLVED-BY-YOU}, registry status \texttt{solved}. Created and updated 14 August 2026. These labels are supplied assessments and do not certify a new result or settle every complete version.
\item Attribution: Ulam AI / UnsolvedMath Contributors, CC BY 4.0. The machine-generated research aid is not peer reviewed and requires independent verification before citation.
\end{itemize}
\end{document}
